\documentclass[11pt]{article}
\usepackage[margin=1in]{geometry}
\usepackage{amsmath,amssymb,amsthm}
\usepackage[colorlinks=true,linkcolor=blue,citecolor=blue,urlcolor=blue]{hyperref}

\newtheorem{theorem}{Theorem}
\newtheorem{lemma}{Lemma}
\newtheorem{remark}{Remark}
\newcommand{\Ric}{\operatorname{Ric}}
\newcommand{\Var}{\operatorname{Var}}

\title{Strict weighted mean-convexity does not control the Neumann gap}
\author{Autonomous partial-result packet for arXiv:1310.2526}
\date{11 August 2026}

\begin{document}
\maketitle

\begin{abstract}
Kolesnikov and Milman ask whether Veysseire's harmonic-mean curvature
spectral-gap estimate has analogues in their Dirichlet and generalized
mean-convex Neumann regimes.  We show that the pure Neumann estimate cannot
hold under generalized mean-convexity alone.  In every dimension at least
three there are smooth bounded Euclidean dumbbells with strictly positive
weighted mean curvature and constant positive Bakry--\'Emery curvature, but
whose weighted Neumann gap tends to zero.  Hence any viable analogue of their
Case~(3) must retain a boundary-trace correction (or impose the stronger
second-fundamental-form condition).  The boundary-corrected Case~(3) and the
Dirichlet question are not settled here.
\end{abstract}

\section{Statement}

On a smooth weighted domain $(\Omega,g,\mu=e^{-V}d\mathrm{vol})$, write
\[
 \Ric_\mu=\Ric_g+\nabla^2V,\qquad
 H_\mu=H-\langle\nabla V,\nu\rangle ,
\]
where $\nu$ is the outer unit normal and $H$ is the sum of the principal
curvatures.  Let $\lambda_1^N(\Omega,\mu)$ be the first nonzero eigenvalue of
the weighted Neumann Laplacian, equivalently the best constant in
\[
 \lambda_1^N\Var_\mu(f)\leq\int_\Omega|\nabla f|^2\,d\mu.
\]

\begin{theorem}[mean-convex dumbbell obstruction]\label{thm:main}
For every $n\geq3$ there are constants $\kappa,c,R>0$ and a family of smooth,
bounded, connected domains $\Omega_\varepsilon\subset B_R(0)\subset\mathbb
R^n$, $0<\varepsilon<\varepsilon_0$, such that for
\[
 d\mu_\varepsilon=e^{-\kappa|x|^2/2}\,dx
\]
one has
\begin{align*}
 \Ric_{\mu_\varepsilon}&=\kappa g_{\mathrm{Eucl}}
       &&\text{on }\Omega_\varepsilon,\\
 H_{\mu_\varepsilon}&\geq c
       &&\text{on }\partial\Omega_\varepsilon,
\end{align*}
but
\[
 \lambda_1^N(\Omega_\varepsilon,\mu_\varepsilon)
 \leq C_n\varepsilon^{\,n-1}\longrightarrow0.
\]
Consequently the Veysseire-form estimate
\begin{equation}\label{eq:false}
 \lambda_1^N\geq
 \left(\frac1{\mu(\Omega)}\int_\Omega\frac1\rho\,d\mu\right)^{-1}
\end{equation}
is false if local convexity is weakened to strict generalized
mean-convexity, even when $\Ric_\mu=\rho g$ with $\rho$ constant.
\end{theorem}

\section{Geometry of the domains}

We use the following standard mean-convex connected-sum construction, and
include the relevant quantitative features.

\begin{lemma}[uniform mean-convex dumbbells]\label{lem:dumbbell}
For every $n\geq3$ there are $R,h_0>0$ and smooth domains
$\Omega_\varepsilon\subset B_R(0)$, symmetric under
$(x_1,x')\mapsto(-x_1,x')$, with the following properties:
\begin{enumerate}
\item $\Omega_\varepsilon$ contains two fixed positive-volume lobes, exchanged
by the symmetry;
\item the lobes are joined by the straight cylinder
$(-1,1)\times B_{\varepsilon}^{\,n-1}$;
\item $H_{\partial\Omega_\varepsilon}\geq h_0$ everywhere.
\end{enumerate}
\end{lemma}

\begin{proof}
Start with two fixed round balls and join them by a rotationally symmetric
tube.  If the tube boundary is written as
\[
 (s,r_\varepsilon(s)\theta),\qquad \theta\in S^{n-2},
\]
then its outward mean curvature is
\begin{equation}\label{eq:Hprofile}
 H=-\frac{r_\varepsilon''}{(1+(r_\varepsilon')^2)^{3/2}}
   +\frac{n-2}{r_\varepsilon\sqrt{1+(r_\varepsilon')^2}}.
\end{equation}
Take $r_\varepsilon=\varepsilon$ on $[-1,1]$.  Immediately outside
this interval, increase the radius while it is still uniformly small, using a
profile with bounded positive second derivative.  The second term in
\eqref{eq:Hprofile} then dominates uniformly.  Continue with a concave
increasing profile, for which both displayed contributions are nonnegative,
and attach it to a fixed spherical lobe.  The inequalities can be made
strict with a margin independent of $\varepsilon$.  Reflect the construction
at $s=0$ and smooth the finitely many joins inside those strict inequalities.
The central cylinder has $H=(n-2)/\varepsilon$, the spherical pieces have
fixed positive mean curvature, and the compact transition pieces retain a
uniform lower bound.  This proves the lemma.
\end{proof}

\section{Weight and Rayleigh quotient}

\begin{proof}[Proof of Theorem \ref{thm:main}]
Apply Lemma \ref{lem:dumbbell}.  Since every domain lies in $B_R(0)$, choose
\[
 \kappa=\frac{h_0}{2R},\qquad V(x)=\frac\kappa2|x|^2.
\]
Euclidean Ricci curvature is zero and $\nabla^2V=\kappa g$, hence
$\Ric_{\mu_\varepsilon}=\kappa g$.  On the boundary,
\[
 H_{\mu_\varepsilon}
 =H-\kappa\langle x,\nu\rangle
 \geq h_0-\kappa|x|
 \geq \frac{h_0}{2}.
\]
Thus the weighted domains are strictly generalized mean-convex, uniformly in
$\varepsilon$.

It remains to estimate their Neumann gap.  Let $f_\varepsilon$ be odd under
the reflection in $x_1$, equal to $-1$ on the left lobe and $1$ on the right
lobe, and interpolate linearly along the central cylinder.  It may be chosen
Lipschitz, with $|\nabla f_\varepsilon|\leq C$ and gradient supported in a
fixed-length portion of the cylinder.  Symmetry of both the domain and the
weight gives $\int f_\varepsilon\,d\mu_\varepsilon=0$.  The two fixed lobes
and the uniform bounds on $V$ give
\[
 \int_{\Omega_\varepsilon}f_\varepsilon^2\,d\mu_\varepsilon\geq c_0>0.
\]
The cylinder has cross-sectional volume $\omega_{n-1}\varepsilon^{n-1}$, so
\[
 \int_{\Omega_\varepsilon}|\nabla f_\varepsilon|^2\,d\mu_\varepsilon
 \leq C_n\varepsilon^{n-1}.
\]
The Neumann Rayleigh principle yields
$\lambda_1^N\leq C_n\varepsilon^{n-1}$.

Finally, in this example one may take $\rho\equiv\kappa$, and therefore
\[
 \left(\frac1{\mu_\varepsilon(\Omega_\varepsilon)}
 \int_{\Omega_\varepsilon}\frac1\rho\,d\mu_\varepsilon\right)^{-1}
 =\kappa.
\]
For small $\varepsilon$, this contradicts \eqref{eq:false}.
\end{proof}

\section{What this does and does not answer}

The locally convex theorem in arXiv:1310.2526 uses
$\mathrm{II}_{\partial M}\geq0$.  Under Neumann boundary conditions, the
Reilly formula contains
\[
 \int_{\partial M}\mathrm{II}_{\partial M}
   (\nabla_{\partial M}u,\nabla_{\partial M}u)\,d\mu_{\partial M},
\]
whose sign is not controlled by the trace $H_\mu$.  The dumbbell construction
is the geometric manifestation of that missing tensor control.

Case~(3) of Kolesnikov--Milman contains an additional boundary-trace variance
weighted by $H_\mu^{-1}$.  The test above has different traces on the two
lobes, so that term does not disappear; the theorem does not refute such a
boundary-corrected Veysseire analogue.  Likewise, a Dirichlet function
vanishes on the entire boundary, and the same low-energy lobe test is
unavailable.  Those two formulations remain open in this packet.

\begin{thebibliography}{9}
\bibitem{KM}
A. V. Kolesnikov and E. Milman,
\emph{Brascamp--Lieb type inequalities on weighted Riemannian manifolds with
boundary}, arXiv:1310.2526.

\bibitem{V}
L. Veysseire,
\emph{Improved spectral gap bounds on positively curved manifolds},
arXiv:1105.6080.
\end{thebibliography}

\end{document}
